% Источники: exam/materials/преподаватель/2026/Моделирование_случайных_величин,_векторов_.pdf; exam/materials/моделирование.pdf, раздел 13.
\section{Методы обращения и суперпозиции, их применение к моделированию скалярных случайных величин.}

\subsection*{Метод обращения (инверсии)}

\textbf{Идея.} Если $W(x)$ — ФР моделируемой СВ $\xi$, то $W(\xi) \sim U(0,1)$. Следовательно, для $r \sim U(0,1)$:
\[
  \xi = W^{-1}(r).
\]

\textbf{Условие применимости:} $W^{-1}$ вычисляется аналитически или табличным способом.

\textbf{Примеры:}
\begin{itemize}
  \item Показательный закон: $W(x)=1-e^{-\lambda x}$ $\Rightarrow$ $\xi = -\frac{1}{\lambda}\ln r$.
  \item Равномерный $U(a,b)$: $\xi = a + (b-a)r$.
  \item Геометрический закон (дискретный): $\xi = \lceil \ln r / \ln(1-p) \rceil$.
\end{itemize}

\subsection*{Метод суперпозиции (кусочного приближения)}

Применяется, когда $W^{-1}$ не имеет аналитического вида. ФР $W(x)$ аппроксимируется \textbf{смесью распределений}:
\[
  W(x) = \sum_{k=1}^K p_k W_k(x), \quad p_k \geqslant 0,\; \sum_k p_k = 1,
\]
где $W_k(x)$ — простые распределения (равномерные, линейные сегменты).

Алгоритм:
\begin{enumerate}
  \item Генерируют $r_1 \sim U(0,1)$; выбирают компоненту $k$ с вероятностью $p_k$.
  \item Генерируют $\xi$ из $W_k$ методом обращения.
\end{enumerate}

\clearpage
